\begin{frame}{실무 결과 2: 라벨러 한 명은 못 믿는다 --- 다수결로 합의}
  \begin{itemize}
    \item 라벨러 한 명의 판단은 \textbf{주관적}일 수 있다 --- 같은
      X-ray를 의사 두 명이 다르게 읽는 일도 흔하다.
    \item 해결: 같은 데이터를 \kb{3명 이상}이 \textbf{독립적으로}
      라벨링 $\to$ \textbf{다수결(majority vote)} 또는 합의로 최종
      라벨 확정.
    \item 라벨러들이 서로 \textbf{얼마나 동의하는지}(inter-annotator
      agreement)를 Cohen's/Fleiss' $\kappa$ 로 수치화 --- $\kappa$ 가
      낮으면 데이터가 아니라 \textbf{라벨 정의 자체가 모호}하다는
      신호.
    \item 예: ImageNet은 Amazon Mechanical Turk 라벨러 다수의 합의로
      각 이미지의 클래스를 확정했다.
  \end{itemize}
  \vskip0.6em
  {\footnotesize\color{ksagray} 출처: Snow et al., ``Cheap and
  Fast---But is it Good? Evaluating Non-Expert Annotations for
  Natural Language Tasks,'' EMNLP 2008; Deng et al., ``ImageNet: A
  Large-Scale Hierarchical Image Database,'' CVPR 2009.}
\end{frame}
